\documentclass[12pt,a4paper]{article}
\usepackage[utf8]{inputenc}
\usepackage[T1]{fontenc}
\usepackage[spanish,es-noquoting]{babel}
\usepackage{geometry}
\usepackage{xcolor}
\usepackage{enumitem}

\geometry{a4paper, margin=1in}

\title{Evaluación de Examen Teórico-Práctico (Grupal)}
\author{Prof. César Rodríguez \\ \small Curso: Seguridad Informática}
\date{\today}

\begin{document}

\maketitle

\section*{Datos de los Estudiantes}
\textbf{Nombres:} Jesmary Rodríguez, Ricardo Rivero, Reinaldo Caraballo \\
\textbf{Grupo:} 3 (Descubrimiento de Hosts) \\
\textbf{Calificación Final:} \textbf{\textcolor{blue}{10/10}}

\section*{Evaluación Detallada}

\begin{itemize}[label=$\bullet$]
    \item \textbf{Pregunta 1 (ICMP en Ping Sweep):} \textbf{2/2 puntos}. Comprenden perfectamente el mecanismo de solicitud (\textit{Echo Request}) y respuesta, mencionando acertadamente el código de tipo (0) para la respuesta (\textit{Echo Reply}) y la implicación de los firewalls.
    
    \item \textbf{Pregunta 2 (Notación CIDR /24):} \textbf{2/2 puntos}. Respuesta correcta y precisa. Detallan tanto la cantidad de direcciones (256) como los límites del rango y su equivalencia exacta en máscara de subred (255.255.255.0).
    
    \item \textbf{Pregunta 3 (Alternativa TCP a ICMP):} \textbf{2/2 puntos}. Identifican perfectamente la técnica de TCP SYN Ping para eludir restricciones de firewall, aportando además un puerto lógico (80) para maximizar el éxito en la prueba.
    
    \item \textbf{Pregunta 4 (Deducción por paquete RST):} \textbf{2/2 puntos}. Demuestran una comprensión sólida de los estados del protocolo TCP; identifican correctamente que un \textit{Reset} (RST) es prueba definitiva de que la máquina está activa, aunque el puerto consultado esté cerrado.
    
    \item \textbf{Pregunta 5 (Librería Python para rangos IP):} \textbf{2/2 puntos}. Nombran correctamente la librería estándar recomendada (\texttt{ipaddress}) y la función constructora específica (\texttt{ip\_network()}) que usarán en su módulo.
\end{itemize}

\vspace{1cm}

\section*{Comentarios Finales}
\noindent El equipo demuestra un dominio sobresaliente y muy preciso de los conceptos de redes, las respuestas del protocolo TCP/IP y las técnicas de descubrimiento. Sus bases teóricas son impecables y están más que preparados para llevar a código la herramienta de escaneo y descubrimiento.

\vspace{0.5cm}
\begin{center}
    \textit{¡Excelente trabajo en equipo!}
\end{center}

\end{document}